% =====================================================================
%  1bat-ud4-12.tex  —  SESSIÓ 12: El disseny de la intervenció: crear l'enquesta
%  (sense preàmbul: s'inclou des de main.tex)
% =====================================================================
\horatitol{Sessió 12 — El disseny de la intervenció: crear l'enquesta}%
{Dissenyar preguntes tancades, clares i neutres: una pregunta mal feta contamina totes les dades que se'n deriven.}

\subsection{Idea clau}

Hem après a \emph{llegir} dades (sessions 8 a 11). Ara comencem a \textbf{recollir-ne
de pròpies}. L'institut vol fer un \textbf{pla de convivència} i necessita dades de
l'alumnat de l'ESO. En grups, dissenyareu una enquesta. El cor de la sessió és la
\textbf{qualitat de les preguntes}: si una pregunta està mal formulada, totes les
dades que en surtin estaran contaminades.

\begin{idea}
\textbf{Com es redacta una bona pregunta d'enquesta}
\begin{itemize}[leftmargin=1.4em, itemsep=2pt]
  \item \textbf{Tancada:} amb opcions de resposta acotades (per poder-les
        \emph{tabular} després). Ex.: «Mai / De vegades / Sovint / Sempre».
  \item \textbf{Clara:} sense ambigüitats; que tothom l'entengui igual.
  \item \textbf{Neutra:} que \textbf{no indueixi} la resposta (sense adjectius ni
        suposicions amagades).
  \item \textbf{Una sola idea:} no preguntar dues coses alhora.
\end{itemize}
\textbf{Regla d'or:} una pregunta esbiaixada \textbf{contamina} totes les dades que
se'n derivin.
\end{idea}

\subsection{Exemples}

\begin{exemple}[pregunta esbiaixada contra pregunta neutra]
Vols saber si l'alumnat se sent segur a l'institut. Compara:
\begin{itemize}[leftmargin=1.4em, itemsep=2pt]
  \item \textbf{Esbiaixada:} «\emph{No creus que l'institut és un lloc segur i
        acollidor?}» — el «no creus que\dots» \textbf{empeny} cap al «sí»; ja suggereix
        la resposta.
  \item \textbf{Neutra:} «\emph{Com de segur et sents a l'institut?}» amb opcions
        «Gens / Poc / Bastant / Molt».
\end{itemize}
La versió neutra deixa que l'alumne respongui lliurement i, a més, dóna opcions
\textbf{tancades} que després es poden comptar i tabular.
\end{exemple}

\begin{exemple}[pregunta que en fa dues alhora]
«\emph{Estàs satisfet amb el menjador i el pati?}» és una pregunta \textbf{doble}: i
si algú està content amb el pati però no amb el menjador? No sabrà què respondre, i la
dada serà confusa. Cal \textbf{separar-la} en dues preguntes:
\begin{itemize}[leftmargin=1.4em, itemsep=1pt]
  \item «Com valores el menjador?» (Molt malament / Malament / Bé / Molt bé)
  \item «Com valores el pati?» (Molt malament / Malament / Bé / Molt bé)
\end{itemize}
\end{exemple}

\subsection{Exercicis resolts}

\begin{exresolt}[corregir una pregunta esbiaixada]
Reescriu aquesta pregunta perquè sigui \textbf{neutra} i \textbf{tancada}:
«\emph{Veritat que passes massa estones amb el mòbil?}»
\solucio
La pregunta original té dos problemes: «\emph{Veritat que\dots}» \textbf{indueix} el
«sí», i «\emph{massa}» ja és un \textbf{judici}. Una versió neutra i tancada:
\begin{center}
«Quantes hores al dia dediques al mòbil fora de l'horari de classe?»\\
\textbf{Opcions:} Menys d'$1$ h \ / \ Entre $1$ i $2$ h \ / \ Entre $2$ i $4$ h \ / \ Més de $4$ h.
\end{center}
Així no jutgem («massa») ni empenyem la resposta, i obtenim categories que es poden
tabular.
\resposta{«Quantes hores al dia dediques al mòbil fora de classe?» amb quatre franges
horàries com a opcions.}
\end{exresolt}

\begin{exresolt}[dissenyar una pregunta tabulable]
Vols saber com dorm l'alumnat. Escriu una pregunta tancada sobre les hores de son i
explica per què les opcions que tries són bones per tabular.
\solucio
Una bona pregunta:
\begin{center}
«Quantes hores dorms, de mitjana, les nits d'escola?»\\
\textbf{Opcions:} Menys de $6$ h \ / \ Entre $6$ i $7$ h \ / \ Entre $7$ i $8$ h \ / \ Més de $8$ h.
\end{center}
Les opcions són \textbf{bones per tabular} perquè: (1) no se solapen (cada resposta cau
en una sola franja); (2) cobreixen tots els casos; (3) són poques i clares, de manera
que després es poden \textbf{comptar} fàcilment i fer-ne percentatges i un gràfic.
\resposta{Pregunta amb quatre franges de son que no se solapen, cobreixen tots els
casos i són fàcils de comptar.}
\end{exresolt}

\subsection{Exercicis per fer}

\begin{exercicis}
\begin{exlist}
  \item Indica si cada pregunta és \textbf{neutra} o \textbf{esbiaixada} i, si cal,
        corregeix-la:
    \begin{enumerate}[label=\alph*), leftmargin=1.6em, topsep=2pt]
      \item «No trobes que hauríem de tenir més activitats?»
      \item «Quantes vegades a la setmana fas esport?»
      \item «T'agrada el nou pati, oi que sí?»
    \end{enumerate}
  \item Converteix aquesta pregunta oberta en una de \textbf{tancada} amb opcions
        tabulables: «Què en penses, de la convivència a l'institut?»
  \item Detecta el problema d'aquesta pregunta i separa-la si cal: «Estàs content amb
        els professors i amb els horaris?»
  \item Dissenya \textbf{tres} preguntes tancades i neutres per a una enquesta sobre
        \textbf{hàbits de salut mental} (per exemple: estrès, son, suport entre
        iguals). Indica les opcions de resposta de cadascuna.
  \item \textbf{(Decisió justificada)} Per què, en una enquesta que després s'haurà de
        \textbf{tabular}, convé que les preguntes siguin \textbf{tancades} i no obertes?
        Quin inconvenient tindria fer-les totes obertes?
  \item \textbf{(Raonament)} Explica, amb les teves paraules i un exemple, la frase
        «una pregunta mal formulada contamina totes les dades que se'n deriven». Per
        què és tan important per a un tècnic que recull dades reals?
\end{exlist}
\end{exercicis}

% ---------------------------------------------------------------------
\begin{idea}
\textbf{Revisió creuada} — abans de tancar l'enquesta, intercanvieu-la amb un altre
grup i comproveu, pregunta a pregunta:
\begin{center}
\renewcommand{\arraystretch}{1.4}
\begin{tabular}{p{8.4cm} c c}
\toprule
\textbf{Comprovació} & \textbf{Sí} & \textbf{No} \\
\midrule
És tancada (té opcions acotades)? & \rule{0.7cm}{0pt} & \\
És clara (tothom l'entendria igual)? & & \\
És neutra (no indueix la resposta)? & & \\
Pregunta una sola cosa? & & \\
Les opcions no se solapen i cobreixen tots els casos? & & \\
\bottomrule
\end{tabular}
\end{center}
\end{idea}

% ---------------------------------------------------------------------
\fullsolucions{Sessió 12}
\begin{solucions}
\begin{sollist}
  \item (a) \textbf{Esbiaixada} («No trobes que\dots» indueix el sí). Neutra: «Com
        valores el nombre d'activitats actual?» (Insuficient / Adequat / Excessiu).
        \quad (b) \textbf{Neutra} i tancable (es poden donar franges: $0$, $1$–$2$,
        $3$–$4$, $5$ o més). \quad (c) \textbf{Esbiaixada} («oi que sí?» empeny el sí).
        Neutra: «Com valores el nou pati?» (Molt malament / Malament / Bé / Molt bé).
  \item Exemple: «Com valores la convivència a l'institut?» amb opcions \textbf{Molt
        dolenta / Dolenta / Bona / Molt bona}.
  \item És una pregunta \textbf{doble} (professors i horaris alhora). Cal separar-la:
        «Com valores els professors?» i «Com valores els horaris?», cadascuna amb les
        seves opcions.
  \item Resposta oberta. Exemples: «Amb quina freqüència et sents estressat per
        l'institut?» (Mai / De vegades / Sovint / Sempre); «Quantes hores dorms les
        nits d'escola?» (Menys de $6$ / $6$–$7$ / $7$–$8$ / Més de $8$); «Si tens un
        problema, tens algú de confiança a qui explicar-l'hi?» (Mai / De vegades /
        Gairebé sempre / Sempre).
  \item Perquè les preguntes \textbf{tancades} donen respostes que es poden
        \textbf{comptar} directament (freqüències, percentatges, gràfics). Si totes
        fossin obertes, cada resposta seria un text diferent, dificilíssim de comptar i
        de comparar: no es podrien tabular bé.
  \item Resposta oberta. Idea: si la pregunta empeny cap a una resposta o és ambigua,
        les dades recollides ja no reflecteixen la realitat, sinó com estava feta la
        pregunta; totes les conclusions posteriors (gràfics, percentatges, decisions)
        seran errònies. Per a un tècnic que recull dades reals, això pot portar a
        dissenyar una intervenció basada en informació falsa.
\end{sollist}
\end{solucions}
